\CNsection{1}{Why Scheduling Matters}

In LTE, the radio channel is a \textbf{shared resource}. Multiple users compete for limited time-frequency resources on the same carrier. The scheduler is the central intelligence at the eNodeB that decides:

\begin{itemize}
\tightlist
\item
  \textbf{Who} gets to transmit/\allowbreak{}receive
\item
  \textbf{When} (which subframe)
\item
  \textbf{Where} (which resource blocks)
\item
  \textbf{How} (which modulation and coding scheme)
\end{itemize}

Without scheduling, users would collide, resources would be wasted, and QoS guarantees would be impossible. The scheduler must balance:

\begin{itemize}
\tightlist
\item
  \textbf{Throughput maximization} — use resources efficiently
\item
  \textbf{Fairness} — don’t starve cell-edge users
\item
  \textbf{Latency} — meet delay budgets for real-time services
\item
  \textbf{QoS} — honor GBR (Guaranteed Bit Rate) bearers
\end{itemize}

\Needspace{22\baselineskip}

\CNsection{2}{The LTE Resource Grid}

The resource grid is the fundamental 2D structure that organizes all LTE transmissions in time and frequency.

\CNsubsection{}{Time Domain Hierarchy}

{\def\LTcaptype{none} % do not increment counter
\Needspace{11\baselineskip}
\begin{longtable}[c]{@{}ccc@{}}
\toprule\noalign{}
\rowcolor{CNink}\CNth{Structure} & \CNth{Duration} & \CNth{Composition} \\
\CNheadrulesep
\endhead
\bottomrule\noalign{}
\endlastfoot
Radio Frame & 10~ms & 10 subframes \\
Subframe & 1~ms & 2 slots \\
Slot & 0.5~ms & 7 OFDM symbols (normal CP) \\
OFDM Symbol & \textasciitilde{}71.4 \ensuremath{\mu}s & 1 symbol period + cyclic prefix \\
\end{longtable}
}

\Needspace{13\baselineskip}

\CNsubsection{}{Frequency Domain Hierarchy}

{\def\LTcaptype{none} % do not increment counter
\Needspace{9\baselineskip}
\begin{longtable}[c]{@{}ccc@{}}
\toprule\noalign{}
\rowcolor{CNink}\CNth{Structure} & \CNth{Bandwidth} & \CNth{Composition} \\
\CNheadrulesep
\endhead
\bottomrule\noalign{}
\endlastfoot
Subcarrier & 15~kHz & Smallest frequency unit \\
Resource Block (RB) & 180~kHz & 12 subcarriers \\
System Bandwidth & 1.4–20~MHz & 6–100~RBs \\
\end{longtable}
}

\CNsubsection{}{Key Definitions}

\begin{itemize}
\tightlist
\item
  \textbf{Resource Element (RE):} 1 subcarrier \ensuremath{\times} 1 OFDM symbol — the smallest allocatable unit
\item
  \textbf{Resource Block (RB):} 12 subcarriers \ensuremath{\times} 7 symbols = \textbf{84~REs} (one slot)
\item
  \textbf{Resource Block Pair:} 12 subcarriers \ensuremath{\times} 14 symbols = \textbf{168~REs} (one subframe /\allowbreak{} TTI)
\end{itemize}

\CNkeepfig{m07-frame-hierarchy}{3}

\CNsubsection{}{Resource Grid Visualization}

\CNfigure{m07-frame-hierarchy}{LTE time structure: frame, subframe, slot and OFDM symbols}

\CNfigure{m07-rb-pair}{Resource elements of one resource block (12 subcarriers \ensuremath{\times} 7 symbols)}

\Needspace{19\baselineskip}

\CNsection{3}{Scheduling Granularity}

{\def\LTcaptype{none} % do not increment counter
\Needspace{13\baselineskip}
\begin{longtable}[c]{@{}
  >{\raggedright\arraybackslash\hspace{0pt}}m{(\linewidth - 2\tabcolsep) * \real{0.6111}}
  >{\raggedright\arraybackslash\hspace{0pt}}m{(\linewidth - 2\tabcolsep) * \real{0.3889}}@{}}
\toprule\noalign{}
\rowcolor{CNink}\CNth{Parameter} & \CNth{Value} \\
\CNheadrulesep
\endhead
\bottomrule\noalign{}
\endlastfoot
\textbf{Time granularity} & 1 subframe = 1~ms (Transmission Time Interval, TTI) \\
\textbf{Frequency granularity} & 1 Resource Block = 180~kHz (12 subcarriers) \\
\textbf{Minimum allocation} & 1 RB pair per TTI per user \\
\textbf{Maximum allocation} & All RBs in a TTI to one user \\
\textbf{Scheduling decision rate} & 1000 decisions/\allowbreak{}second \\
\end{longtable}
}

The scheduler makes a fresh decision every TTI (1~ms), assigning RBs to UEs based on channel conditions, buffer status, QoS requirements, and fairness metrics.

\Needspace{14\baselineskip}

\CNsection{4}{Downlink Scheduling}

\CNchained\CNsubsection{}{Overview}

The eNodeB scheduler controls all downlink resource assignments. It informs UEs of their allocations via \textbf{DCI (Downlink Control Information)} messages carried on the \textbf{PDCCH (Physical Downlink Control Channel)}.

\CNsubsection{}{DCI Contents}

A DCI message tells the UE:

\begin{itemize}
\tightlist
\item
  Which RBs are allocated (resource assignment)
\item
  Which MCS to use for decoding
\item
  HARQ process number and redundancy version
\item
  Power control commands
\item
  MIMO-related parameters (precoding, layers)
\end{itemize}

\Needspace{13\baselineskip}

\CNsubsection{}{Resource Allocation Types}

{\def\LTcaptype{none} % do not increment counter
\Needspace{9\baselineskip}
\begin{longtable}[c]{@{}
  >{\raggedright\arraybackslash\hspace{0pt}}m{(\linewidth - 6\tabcolsep) * \real{0.1622}}
  >{\raggedright\arraybackslash\hspace{0pt}}m{(\linewidth - 6\tabcolsep) * \real{0.2162}}
  >{\raggedright\arraybackslash\hspace{0pt}}m{(\linewidth - 6\tabcolsep) * \real{0.3514}}
  >{\raggedright\arraybackslash\hspace{0pt}}m{(\linewidth - 6\tabcolsep) * \real{0.2703}}@{}}
\toprule\noalign{}
\rowcolor{CNink}\CNth{Type} & \CNth{Method} & \CNth{Flexibility} & \CNth{Overhead} \\
\CNheadrulesep
\endhead
\bottomrule\noalign{}
\endlastfoot
\textbf{Type 0} & Bitmap of RB Groups (RBGs) & Distributed, non-contiguous & Medium \\
\textbf{Type 1} & Subset + shift + bitmap & Distributed within subset & Higher \\
\textbf{Type 2} & Start RB + length (RIV) & Contiguous only & Lowest \\
\end{longtable}
}

\Needspace{17\baselineskip}

\textbf{RBG size} depends on system bandwidth:

{\def\LTcaptype{none} % do not increment counter
\Needspace{14\baselineskip}
\begin{longtable}[c]{@{}cc@{}}
\toprule\noalign{}
\rowcolor{CNink}\CNth{System BW (RBs)} & \CNth{RBG Size (RBs)} \\
\CNheadrulesep
\endhead
\bottomrule\noalign{}
\endlastfoot
6 & 1 \\
15 & 2 \\
25 & 2 \\
50 & 3 \\
75 & 4 \\
100 & 4 \\
\end{longtable}
}

\CNsubsection{}{MCS Selection}

The scheduler selects MCS based on:

\begin{enumerate}
\def\labelenumi{\arabic{enumi}.}
\tightlist
\item
  UE reports \textbf{CQI} (Channel Quality Indicator)
\item
  Scheduler maps CQI \ensuremath{\to} appropriate MCS index
\item
  Applies outer-loop adjustments based on ACK/\allowbreak{}NACK history
\item
  Encodes DCI with selected MCS for the UE to decode PDSCH
\end{enumerate}

\Needspace{14\baselineskip}

\CNsection{5}{Uplink Scheduling}

\CNchained\CNsubsection{}{Key Differences from Downlink}

Unlike downlink, uplink scheduling must account for:

\begin{itemize}
\tightlist
\item
  \textbf{SC-FDMA constraint:} UEs must be allocated \textbf{contiguous} RBs (single-carrier property)
\item
  \textbf{UE power limitations:} Distance from eNodeB affects achievable rate
\item
  \textbf{Buffer awareness:} eNodeB doesn’t directly know UE buffer status
\end{itemize}

\CNsubsection{}{Scheduling Grant}

The eNodeB sends an \textbf{uplink grant} via DCI (Format 0/\allowbreak{}4) on PDCCH, containing:

\begin{itemize}
\tightlist
\item
  RB allocation (contiguous)
\item
  MCS index
\item
  Power control command
\item
  HARQ information
\end{itemize}

\Needspace{15\baselineskip}

\CNsubsection{}{Buffer Status Reports (BSR)}

UEs inform the eNodeB about queued data via BSR:

{\def\LTcaptype{none} % do not increment counter
\Needspace{9\baselineskip}
\begin{longtable}[c]{@{}
  >{\raggedright\arraybackslash\hspace{0pt}}m{(\linewidth - 2\tabcolsep) * \real{0.5263}}
  >{\raggedright\arraybackslash\hspace{0pt}}m{(\linewidth - 2\tabcolsep) * \real{0.4737}}@{}}
\toprule\noalign{}
\rowcolor{CNink}\CNth{BSR Type} & \CNth{Trigger} \\
\CNheadrulesep
\endhead
\bottomrule\noalign{}
\endlastfoot
\textbf{Regular BSR} & New data arrives in a logical channel group with higher priority than currently reported \\
\textbf{Periodic BSR} & Timer-based periodic reporting \\
\textbf{Padding BSR} & Sent when UL grant has spare capacity \\
\end{longtable}
}

BSR reports data volume per \textbf{Logical Channel Group (LCG)}, allowing priority-aware scheduling.

\CNsubsection{}{Power Headroom Report (PHR)}

PHR tells the eNodeB how much \textbf{additional transmit power} the UE has available:

\CNdisplay{PH = P_{MAX} - P_{PUSCH}}

\begin{itemize}
\tightlist
\item
  If PH \textgreater{} 0: UE can support higher MCS or more RBs
\item
  If PH \ensuremath{\le} 0: UE is power-limited, scheduler should reduce allocation
\end{itemize}

\CNkeepfig{m07-scfdma-alloc}{3}

\CNsubsection{}{SC-FDMA Contiguous Allocation}

\CNfigure{m07-scfdma-alloc}{SC-FDMA uplink allocations must be contiguous in frequency}

\Needspace{14\baselineskip}

\CNsection{6}{Channel Quality Indicator (CQI)}

\CNchained\CNsubsection{}{What CQI Represents}

CQI is a \textbf{4-bit index} (1–15) reported by the UE indicating the highest MCS that can be received with a transport block error rate (BLER) \ensuremath{\le} \textbf{10\%}.

CQI = 0 means ``out of range'' (channel too poor for any transmission).

\Needspace{26\baselineskip}

\CNsubsection{}{CQI Table (3GPP TS 36.213, Table 7.2.3-1)}

{\def\LTcaptype{none} % do not increment counter
\Needspace{22\baselineskip}
\begin{longtable}[c]{@{}
  >{\raggedright\arraybackslash\hspace{0pt}}m{(\linewidth - 6\tabcolsep) * \real{0.1528}}
  >{\raggedright\arraybackslash\hspace{0pt}}m{(\linewidth - 6\tabcolsep) * \real{0.1667}}
  >{\raggedright\arraybackslash\hspace{0pt}}m{(\linewidth - 6\tabcolsep) * \real{0.2639}}
  >{\raggedright\arraybackslash\hspace{0pt}}m{(\linewidth - 6\tabcolsep) * \real{0.4167}}@{}}
\toprule\noalign{}
\rowcolor{CNink}\CNth{CQI Index} & \CNth{Modulation} & \CNth{Code Rate \ensuremath{\times} 1024} & \CNth{Spectral Efficiency (bps/\allowbreak{}Hz)} \\
\CNheadrulesep
\endhead
\bottomrule\noalign{}
\endlastfoot
1 & QPSK & 78 & 0.1523 \\
2 & QPSK & 120 & 0.2344 \\
3 & QPSK & 193 & 0.3770 \\
4 & QPSK & 308 & 0.6016 \\
5 & QPSK & 449 & 0.8770 \\
6 & QPSK & 602 & 1.1758 \\
7 & 16QAM & 378 & 1.4766 \\
8 & 16QAM & 490 & 1.9141 \\
9 & 16QAM & 616 & 2.4063 \\
10 & 64QAM & 466 & 2.7305 \\
11 & 64QAM & 567 & 3.3223 \\
12 & 64QAM & 666 & 3.9023 \\
13 & 64QAM & 772 & 4.5234 \\
14 & 64QAM & 873 & 5.1152 \\
15 & 64QAM & 948 & 5.5547 \\
\end{longtable}
}

\Needspace{13\baselineskip}

\CNsubsection{}{Wideband vs Subband CQI}

{\def\LTcaptype{none} % do not increment counter
\Needspace{9\baselineskip}
\begin{longtable}[c]{@{}
  >{\raggedright\arraybackslash\hspace{0pt}}m{(\linewidth - 4\tabcolsep) * \real{0.3947}}
  >{\raggedright\arraybackslash\hspace{0pt}}m{(\linewidth - 4\tabcolsep) * \real{0.3421}}
  >{\raggedright\arraybackslash\hspace{0pt}}m{(\linewidth - 4\tabcolsep) * \real{0.2632}}@{}}
\toprule\noalign{}
\rowcolor{CNink}\CNth{Reporting Mode} & \CNth{Description} & \CNth{Use Case} \\
\CNheadrulesep
\endhead
\bottomrule\noalign{}
\endlastfoot
\textbf{Wideband CQI} & Single CQI for entire bandwidth & Low overhead, coarse adaptation \\
\textbf{Subband CQI} & CQI per group of RBs & Frequency-selective scheduling, higher gain \\
\textbf{UE-selected subband} & UE reports best-M subbands & Compromise between overhead and gain \\
\end{longtable}
}

Subband CQI enables \textbf{frequency-selective scheduling}, where the scheduler assigns each UE the RBs where their channel is best.

\Needspace{14\baselineskip}

\CNsection{7}{Link Adaptation}

\CNchained\CNsubsection{}{Adaptive Modulation and Coding (AMC)}

AMC dynamically adjusts the modulation order and code rate to match instantaneous channel conditions:

\begin{itemize}
\tightlist
\item
  \textbf{Good channel} \ensuremath{\to} Higher-order modulation + higher code rate \ensuremath{\to} More bits/\allowbreak{}RE
\item
  \textbf{Poor channel} \ensuremath{\to} Lower-order modulation + lower code rate \ensuremath{\to} Fewer bits/\allowbreak{}RE but robust
\end{itemize}

\Needspace{15\baselineskip}

\CNsubsection{}{Modulation Schemes}

{\def\LTcaptype{none} % do not increment counter
\Needspace{11\baselineskip}
\begin{longtable}[c]{@{}cccc@{}}
\toprule\noalign{}
\rowcolor{CNink}\CNth{Modulation} & \CNth{Bits per Symbol} & \CNth{Required SNR (approx.)} & \CNth{Usage} \\
\CNheadrulesep
\endhead
\bottomrule\noalign{}
\endlastfoot
QPSK & 2 & \textless{} 10~dB & Cell-edge users \\
16QAM & 4 & 10–17~dB & Mid-range users \\
64QAM & 6 & 17–25~dB & Near-cell users \\
256QAM & 8 & \textgreater{} 25~dB & Cat-11+ (Rel-12), best conditions \\
\end{longtable}
}

\Needspace{26\baselineskip}

\CNsubsection{}{MCS Index Table (3GPP TS 36.213, Table 7.1.7.1-1)}

{\def\LTcaptype{none} % do not increment counter
\Needspace{22\baselineskip}
\begin{longtable}[c]{@{}cccc@{}}
\toprule\noalign{}
\rowcolor{CNink}\CNth{MCS Index} & \CNth{Modulation Order} & \CNth{TBS Index} & \CNth{Notes} \\
\CNheadrulesep
\endhead
\bottomrule\noalign{}
\endlastfoot
0 & 2 (QPSK) & 0 & Lowest rate \\
1 & 2 (QPSK) & 1 & \\
2 & 2 (QPSK) & 2 & \\
3 & 2 (QPSK) & 3 & \\
4 & 2 (QPSK) & 4 & \\
5 & 2 (QPSK) & 5 & \\
6 & 2 (QPSK) & 6 & \\
7 & 2 (QPSK) & 7 & \\
8 & 2 (QPSK) & 8 & \\
9 & 2 (QPSK) & 9 & \\
10 & 4 (16QAM) & 9 & Modulation transition \\
11 & 4 (16QAM) & 10 & \\
12 & 4 (16QAM) & 11 & \\
13 & 4 (16QAM) & 12 & \\
14 & 4 (16QAM) & 13 & \\
15 & 4 (16QAM) & 14 & \\
16 & 4 (16QAM) & 15 & \\
17 & 6 (64QAM) & 15 & Modulation transition \\
18 & 6 (64QAM) & 16 & \\
19 & 6 (64QAM) & 17 & \\
20 & 6 (64QAM) & 18 & \\
21 & 6 (64QAM) & 19 & \\
22 & 6 (64QAM) & 20 & \\
23 & 6 (64QAM) & 21 & \\
24 & 6 (64QAM) & 22 & \\
25 & 6 (64QAM) & 23 & \\
26 & 6 (64QAM) & 24 & \\
27 & 6 (64QAM) & 25 & \\
28 & 6 (64QAM) & 26 & Highest rate \\
29 & 2 (QPSK) & — & Reserved (retx) \\
30 & 4 (16QAM) & — & Reserved (retx) \\
31 & 6 (64QAM) & — & Reserved (retx) \\
\end{longtable}
}

\CNsubsection{}{Outer Loop Link Adaptation (OLLA)}

The inner loop (CQI \ensuremath{\to} MCS) can be inaccurate due to:

\begin{itemize}
\tightlist
\item
  CQI reporting delay
\item
  UE measurement errors
\item
  Interference variability
\end{itemize}

\textbf{OLLA} adjusts the effective CQI with an offset:

\CNdisplay{MCS_{effective} = f(CQI + \Delta_{OLLA})}

The offset is updated based on HARQ feedback:

\begin{itemize}
\tightlist
\item
  \textbf{ACK received:} $\Delta_{\mathrm{OLLA}} \mathrel{+}= \mathrm{step\_up}$ (e.g., +0.1~dB) \ensuremath{\to} try higher MCS
\item
  \textbf{NACK received:} $\Delta_{\mathrm{OLLA}} \mathrel{-}= \mathrm{step\_down}$ (e.g., -1.0~dB) \ensuremath{\to} fall back to lower MCS
\end{itemize}

Target: maintain \textbf{10\% BLER} at first transmission (ratio $\mathrm{step\_down}/\mathrm{step\_up} \approx 9{:}1$).

\Needspace{14\baselineskip}

\CNsection{8}{HARQ (Hybrid Automatic Repeat Request)}

\CNchained\CNsubsection{}{Overview}

HARQ combines \textbf{FEC (Forward Error Correction)} with \textbf{ARQ (retransmission)} for reliable delivery with low latency.

\Needspace{17\baselineskip}

\CNsubsection{}{LTE HARQ Parameters}

{\def\LTcaptype{none} % do not increment counter
\Needspace{13\baselineskip}
\begin{longtable}[c]{@{}ccc@{}}
\toprule\noalign{}
\rowcolor{CNink}\CNth{Parameter} & \CNth{Downlink} & \CNth{Uplink} \\
\CNheadrulesep
\endhead
\bottomrule\noalign{}
\endlastfoot
Type & Asynchronous & Synchronous \\
\# Processes & 8 & 8 \\
Timing & Flexible retx subframe & Fixed: retx at n+8 \\
Feedback & ACK/\allowbreak{}NACK on PUCCH/\allowbreak{}PUSCH & PHICH (ACK/\allowbreak{}NACK) \\
RTT & 8~ms & 8~ms \\
\end{longtable}
}

\CNkeepfig{m07-harq-processes}{5}

\CNsubsection{}{Stop-and-Wait with 8 Processes}

Each process handles one transport block at a time (stop-and-wait). With \textbf{8 parallel processes}, the pipeline stays full:

\CNfigure{m07-harq-processes}{Eight parallel stop-and-wait HARQ processes}

\textbf{RTT = 8~ms:} UE receives at time $t$ \ensuremath{\to} decodes \ensuremath{\to} sends ACK/\allowbreak{}NACK at $t+4$ \ensuremath{\to} eNB receives \ensuremath{\to} retransmits at $t+8$.

\Needspace{12\baselineskip}

\CNsubsection{}{Combining Methods}

{\def\LTcaptype{none} % do not increment counter
\Needspace{8\baselineskip}
\begin{longtable}[c]{@{}
  >{\raggedright\arraybackslash\hspace{0pt}}m{(\linewidth - 4\tabcolsep) * \real{0.2500}}
  >{\raggedright\arraybackslash\hspace{0pt}}m{(\linewidth - 4\tabcolsep) * \real{0.4062}}
  >{\raggedright\arraybackslash\hspace{0pt}}m{(\linewidth - 4\tabcolsep) * \real{0.3438}}@{}}
\toprule\noalign{}
\rowcolor{CNink}\CNth{Method} & \CNth{Description} & \CNth{Advantage} \\
\CNheadrulesep
\endhead
\bottomrule\noalign{}
\endlastfoot
\textbf{Chase Combining (CC)} & Retransmit identical bits, soft-combine at receiver & Simple, \textasciitilde{}3~dB gain \\
\textbf{Incremental Redundancy (IR)} & Retransmit different redundancy version (new parity bits) & Higher coding gain, better performance \\
\end{longtable}
}

LTE uses \textbf{incremental redundancy} by default with 4 redundancy versions (RV 0, 2, 3, 1).

\Needspace{14\baselineskip}

\Needspace{24\baselineskip}

\CNsection{9}{Scheduling Algorithms}

\CNchained\CNsubsection{}{Round Robin (RR)}

\textbf{Principle:} Allocate equal resources to all users in turn.

\CNdisplay{\text{User}_{selected}(t) = (t \mod N_{users})}

{\def\LTcaptype{none} % do not increment counter
\Needspace{11\baselineskip}
\begin{longtable}[c]{@{}
  >{\raggedright\arraybackslash\hspace{0pt}}m{(\linewidth - 2\tabcolsep) * \real{0.5000}}
  >{\raggedright\arraybackslash\hspace{0pt}}m{(\linewidth - 2\tabcolsep) * \real{0.5000}}@{}}
\toprule\noalign{}
\rowcolor{CNink}\CNth{Pros} & \CNth{Cons} \\
\CNheadrulesep
\endhead
\bottomrule\noalign{}
\endlastfoot
Simple to implement & Ignores channel conditions \\
Perfectly fair in resources & Unfair in throughput (edge users get less) \\
No feedback required & Poor spectral efficiency \\
Low computational cost & No QoS differentiation \\
\end{longtable}
}

\CNsubsection{}{Proportional Fair (PF)}

\textbf{Principle:} Maximize the ratio of instantaneous rate to average throughput.

\CNdisplay{\text{User}_{selected} = \arg\max_k \frac{R_k^{inst}(t)}{\bar{R}_k(t)}}

Where:

\begin{itemize}
\tightlist
\item
  $R_k^{inst}(t)$ = instantaneous achievable rate for user k (from CQI)
\item
  $\bar{R}_k(t)$ = exponential moving average of user k’s throughput
\end{itemize}

\Needspace{16\baselineskip}

Average update:\CNdisplay{\bar{R}_k(t+1) = (1-\frac{1}{t_c})\bar{R}_k(t) + \frac{1}{t_c} \cdot r_k(t)}

where $t_c$ is the averaging window (typically 100–1000 TTIs).

{\def\LTcaptype{none} % do not increment counter
\Needspace{11\baselineskip}
\begin{longtable}[c]{@{}cc@{}}
\toprule\noalign{}
\rowcolor{CNink}\CNth{Pros} & \CNth{Cons} \\
\CNheadrulesep
\endhead
\bottomrule\noalign{}
\endlastfoot
Exploits multi-user diversity & More complex \\
Good throughput-fairness balance & Requires CQI feedback \\
Schedules users on their peaks & Tuning tc affects behavior \\
Maximizes sum log(throughput) & Not optimal for strict delay \\
\end{longtable}
}

\Needspace{17\baselineskip}

\CNsubsection{}{Maximum C/\allowbreak{}I (Max Throughput)}

\textbf{Principle:} Always schedule the user with the best channel.

\CNdisplay{\text{User}_{selected} = \arg\max_k R_k^{inst}(t)}

{\def\LTcaptype{none} % do not increment counter
\Needspace{9\baselineskip}
\begin{longtable}[c]{@{}cc@{}}
\toprule\noalign{}
\rowcolor{CNink}\CNth{Pros} & \CNth{Cons} \\
\CNheadrulesep
\endhead
\bottomrule\noalign{}
\endlastfoot
Maximizes cell throughput & Starves cell-edge users \\
Best spectral efficiency & Extremely unfair \\
Simple decision & Violates QoS for weak users \\
\end{longtable}
}

\CNsubsection{}{QoS-Aware Scheduling}

\textbf{Principle:} Guarantee minimum rates for GBR bearers, prioritize by QCI.

Steps:

\begin{enumerate}
\def\labelenumi{\arabic{enumi}.}
\tightlist
\item
  \textbf{Satisfy GBR bearers first} — allocate minimum required RBs
\item
  \textbf{Priority ordering} — rank by QCI priority level
\item
  \textbf{Remaining resources} — distribute via PF among non-GBR bearers
\item
  \textbf{Delay budget check} — boost priority for packets approaching deadline
\end{enumerate}

\Needspace{14\baselineskip}

\CNkeepfig{m07-trilemma}{8}

\CNsection{10}{Engineering Trade-offs}

\CNchained\CNsubsection{}{The Scheduling Trilemma}

\CNfigure{m07-trilemma}{The scheduling trilemma and where each algorithm sits}

{\def\LTcaptype{none} % do not increment counter
\Needspace{11\baselineskip}
\begin{longtable}[c]{@{}ccccc@{}}
\toprule\noalign{}
\rowcolor{CNink}\CNth{Algorithm} & \CNth{Throughput} & \CNth{Fairness} & \CNth{Latency} & \CNth{Complexity} \\
\CNheadrulesep
\endhead
\bottomrule\noalign{}
\endlastfoot
Round Robin & Low & High (resource) & Medium & Very Low \\
Proportional Fair & High & Medium & Medium & Medium \\
Max C/\allowbreak{}I & Highest & Very Low & Low (for best users) & Low \\
QoS-Aware & Medium-High & Configurable & Guaranteed & High \\
\end{longtable}
}

\Needspace{17\baselineskip}

\CNsubsection{}{Design Considerations}

{\def\LTcaptype{none} % do not increment counter
\Needspace{13\baselineskip}
\begin{longtable}[c]{@{}
  >{\raggedright\arraybackslash\hspace{0pt}}m{(\linewidth - 2\tabcolsep) * \real{0.2759}}
  >{\raggedright\arraybackslash\hspace{0pt}}m{(\linewidth - 2\tabcolsep) * \real{0.7241}}@{}}
\toprule\noalign{}
\rowcolor{CNink}\CNth{Factor} & \CNth{Impact on Scheduling} \\
\CNheadrulesep
\endhead
\bottomrule\noalign{}
\endlastfoot
\textbf{User density} & More users \ensuremath{\to} more multi-user diversity \ensuremath{\to} PF gains increase \\
\textbf{Traffic type} & VoLTE needs strict latency; FTP tolerates delay \\
\textbf{Cell size} & Large cells \ensuremath{\to} big CQI spread \ensuremath{\to} fairness more critical \\
\textbf{Frequency band} & Higher bands \ensuremath{\to} more path loss variation \ensuremath{\to} more scheduling gain \\
\textbf{MIMO mode} & MU-MIMO allows co-scheduling users on same RBs \\
\end{longtable}
}

\Needspace{14\baselineskip}

\CNsection{11}{Numerical Example}

\CNchained\CNsubsection{}{System Parameters}

\begin{itemize}
\tightlist
\item
  \textbf{Bandwidth:} 10~MHz \ensuremath{\to} \textbf{50~RBs} available per TTI
\item
  \textbf{Users:} 5 UEs with different channel conditions
\item
  \textbf{Overhead:} Assume 10 data symbols per subframe (after control/\allowbreak{}reference signals)
\end{itemize}

\Needspace{17\baselineskip}

\CNsubsection{}{User CQI and Spectral Efficiency}

{\def\LTcaptype{none} % do not increment counter
\Needspace{13\baselineskip}
\begin{longtable}[c]{@{}
  >{\raggedright\arraybackslash\hspace{0pt}}m{(\linewidth - 8\tabcolsep) * \real{0.0811}}
  >{\raggedright\arraybackslash\hspace{0pt}}m{(\linewidth - 8\tabcolsep) * \real{0.0676}}
  >{\raggedright\arraybackslash\hspace{0pt}}m{(\linewidth - 8\tabcolsep) * \real{0.1622}}
  >{\raggedright\arraybackslash\hspace{0pt}}m{(\linewidth - 8\tabcolsep) * \real{0.4054}}
  >{\raggedright\arraybackslash\hspace{0pt}}m{(\linewidth - 8\tabcolsep) * \real{0.2838}}@{}}
\toprule\noalign{}
\rowcolor{CNink}\CNth{User} & \CNth{CQI} & \CNth{Modulation} & \CNth{Spectral Efficiency (bps/\allowbreak{}Hz)} & \CNth{Bits per RB per TTI} \\
\CNheadrulesep
\endhead
\bottomrule\noalign{}
\endlastfoot
UE1 & 4 & QPSK & 0.6016 & 0.6016 \ensuremath{\times} 180kHz \ensuremath{\times} 1ms = 108 \\
UE2 & 7 & 16QAM & 1.4766 & 1.4766 \ensuremath{\times} 180kHz \ensuremath{\times} 1ms = 266 \\
UE3 & 10 & 64QAM & 2.7305 & 2.7305 \ensuremath{\times} 180kHz \ensuremath{\times} 1ms = 491 \\
UE4 & 13 & 64QAM & 4.5234 & 4.5234 \ensuremath{\times} 180kHz \ensuremath{\times} 1ms = 814 \\
UE5 & 15 & 64QAM & 5.5547 & 5.5547 \ensuremath{\times} 180kHz \ensuremath{\times} 1ms = 1000 \\
\end{longtable}
}

\begin{CNquote}

\textbf{Note:} Bits per RB per TTI \ensuremath{\approx} Spectral Efficiency \ensuremath{\times} 12 subcarriers \ensuremath{\times} 15~kHz spacing \ensuremath{\times} 1~ms \ensuremath{\approx} SE \ensuremath{\times} 180 (simplified; actual TBS lookup gives discrete values).

\end{CNquote}

\Needspace{20\baselineskip}

\CNsubsection{}{Round Robin Allocation}

Each user gets equal RBs: \textbf{50~RBs \ensuremath{\div} 5 users = 10 RBs/\allowbreak{}user}

{\def\LTcaptype{none} % do not increment counter
\Needspace{14\baselineskip}
\begin{longtable}[c]{@{}cccc@{}}
\toprule\noalign{}
\rowcolor{CNink}\CNth{User} & \CNth{RBs} & \CNth{Throughput per TTI} & \CNth{Throughput (Mbps)} \\
\CNheadrulesep
\endhead
\bottomrule\noalign{}
\endlastfoot
UE1 & 10 & 10 \ensuremath{\times} 108 = 1,080~bits & 1.08 \\
UE2 & 10 & 10 \ensuremath{\times} 266 = 2,660~bits & 2.66 \\
UE3 & 10 & 10 \ensuremath{\times} 491 = 4,910~bits & 4.91 \\
UE4 & 10 & 10 \ensuremath{\times} 814 = 8,140~bits & 8.14 \\
UE5 & 10 & 10 \ensuremath{\times} 1,000 = 10,000~bits & 10.00 \\
\textbf{Total} & \textbf{50} & \textbf{26,790~bits} & \textbf{26.79} \\
\end{longtable}
}

\textbf{Jain’s Fairness Index (resource):} 1.0 (perfectly fair in RBs)\CNbr{}\textbf{Jain’s Fairness Index (throughput):} \textasciitilde{}0.69

\Needspace{20\baselineskip}

\CNsubsection{}{Maximum C/\allowbreak{}I Allocation}

All 50~RBs go to UE5 (best channel):

{\def\LTcaptype{none} % do not increment counter
\Needspace{14\baselineskip}
\begin{longtable}[c]{@{}ccc@{}}
\toprule\noalign{}
\rowcolor{CNink}\CNth{User} & \CNth{RBs} & \CNth{Throughput (Mbps)} \\
\CNheadrulesep
\endhead
\bottomrule\noalign{}
\endlastfoot
UE1 & 0 & 0 \\
UE2 & 0 & 0 \\
UE3 & 0 & 0 \\
UE4 & 0 & 0 \\
UE5 & 50 & 50.00 \\
\textbf{Total} & \textbf{50} & \textbf{50.00} \\
\end{longtable}
}

\textbf{Cell throughput:} 50~Mbps (maximum possible)\CNbr{}\textbf{Jain’s Fairness Index:} 0.20 (worst case for 5 users)

\CNsubsection{}{Proportional Fair Allocation}

Assuming all users start with equal average throughput, PF initially behaves like Max C/\allowbreak{}I. After convergence, it balances by scheduling users on their relative peaks.

\Needspace{19\baselineskip}

\textbf{Steady-state approximation} (PF tends toward equal throughput in log domain):

The PF scheduler approximately equalizes $\log(\bar{R}_k)$, leading to a throughput distribution between RR and Max C/\allowbreak{}I:

{\def\LTcaptype{none} % do not increment counter
\Needspace{14\baselineskip}
\begin{longtable}[c]{@{}cccc@{}}
\toprule\noalign{}
\rowcolor{CNink}\CNth{User} & \CNth{Approx. RBs} & \CNth{Throughput (Mbps)} & \CNth{Metric (\ensuremath{R_{\mathrm{inst}}}/\allowbreak{}\ensuremath{R_{\mathrm{avg}}})} \\
\CNheadrulesep
\endhead
\bottomrule\noalign{}
\endlastfoot
UE1 & 15 & 1.62 & Boosted (low \ensuremath{R_{\mathrm{avg}}}) \\
UE2 & 12 & 3.19 & Moderate \\
UE3 & 10 & 4.91 & Baseline \\
UE4 & 8 & 6.51 & Slightly reduced \\
UE5 & 5 & 5.00 & Reduced (high \ensuremath{R_{\mathrm{avg}}}) \\
\textbf{Total} & \textbf{50} & \textbf{21.23} & — \\
\end{longtable}
}

\begin{CNquote}

\textbf{Note:} Exact PF allocation depends on fading dynamics and tc window. In flat-fading (no frequency selectivity), PF gives more resources to weaker users over time. With frequency-selective fading, all users benefit from being scheduled on their best subbands.

\end{CNquote}

\Needspace{18\baselineskip}

\CNsubsection{}{Comparison Summary}

{\def\LTcaptype{none} % do not increment counter
\Needspace{14\baselineskip}
\begin{longtable}[c]{@{}cccc@{}}
\toprule\noalign{}
\rowcolor{CNink}\CNth{Metric} & \CNth{Round Robin} & \CNth{Max C/\allowbreak{}I} & \CNth{Proportional Fair} \\
\CNheadrulesep
\endhead
\bottomrule\noalign{}
\endlastfoot
Cell Throughput (Mbps) & 26.79 & 50.00 & \textasciitilde{}21–35 \\
UE1 Throughput (Mbps) & 1.08 & 0 & \textasciitilde{}1.6 \\
UE5 Throughput (Mbps) & 10.00 & 50.00 & \textasciitilde{}5.0 \\
Fairness (Jain’s, resource) & 1.0 & 0.20 & \textasciitilde{}0.75 \\
Spectral Efficiency & Low & Maximum & High \\
Starvation? & No & Yes (4 users) & No \\
\end{longtable}
}

\CNboxsection{Key Takeaways}

\begin{CNbox}[unbreakable]{sum}{Key Takeaways}

\begin{enumerate}
\def\labelenumi{\arabic{enumi}.}
\tightlist
\item
  \textbf{The resource grid} is a time-frequency matrix — scheduling operates at RB granularity every 1~ms TTI
\item
  \textbf{CQI feedback} is the scheduler’s eyes — it drives MCS selection and resource allocation
\item
  \textbf{HARQ} provides reliability with 8 parallel processes and 8~ms RTT
\item
  \textbf{Link adaptation (AMC + OLLA)} continuously matches rate to channel conditions
\item
  \textbf{No single algorithm wins} — the choice depends on deployment scenario, traffic mix, and operator priorities
\item
  \textbf{Proportional Fair} is the de-facto standard in commercial networks as it provides a good throughput-fairness compromise
\item
  \textbf{SC-FDMA} constrains uplink to contiguous allocation, limiting scheduling flexibility compared to downlink OFDMA
\end{enumerate}

\end{CNbox}

\CNsection{}{References}

\begin{itemize}
\tightlist
\item
  3GPP TS 36.213: Physical layer procedures (CQI/\allowbreak{}MCS tables)
\item
  3GPP TS 36.321: MAC protocol (BSR, PHR, HARQ)
\item
  3GPP TS 36.212: Multiplexing and channel coding (TBS tables)
\item
  Sesia, Toufik, Baker: ``LTE – The UMTS Long Term Evolution'' (Wiley)
\end{itemize}
